%! Confidence Intervals for the Difference of Two Proportions — Unit 6, Lesson 6.8 — Homework
\documentclass[10pt]{article}
\usepackage{apstats-article}
\usepackage{apstats-boxes}
\newcommand{\TallMath}[1]{$\displaystyle #1\rule[-1.4em]{0pt}{3.2em}$}

\begin{document}

\pageheader{Unit 6, Lesson 6.8}{Homework}
\namedateperiod

\vspace{0.15in}

\begin{enumerate}

\item \textbf{Exercise Frequency.} A fitness researcher collected independent random samples
from two cities to compare the proportions of adults who exercise at least 3 times per week.
In City X, 168 of 280 adults met this criterion. In City Y, 126 of 240 adults met it.

\begin{enumerate}[label=\alph*.]
  \item Define $p_X$ and $p_Y$ and compute the point estimate $\hat p_X - \hat p_Y$.
        \writeline
  \item Verify all conditions (Random, Large Counts, 10\%).
        \writelines{3}
  \item Construct a 95\% confidence interval for $p_X - p_Y$.
        \writelines{3}
  \item Interpret the interval in context.
        \writelines{2}
\end{enumerate}

\item \textbf{Seat Belt Use.} A highway safety agency surveyed drivers in two states.
In State 1, 374 of 440 drivers reported always wearing a seat belt.
In State 2, 306 of 390 drivers reported the same.

\begin{enumerate}[label=\alph*.]
  \item Verify the Large Counts condition for both groups.
        \writelines{2}
  \item Construct a 99\% confidence interval for $p_1 - p_2$.
        \writelines{3}
  \item Does your interval provide evidence of a difference in seat belt use between the two states? Explain.
        \writelines{2}
\end{enumerate}

\end{enumerate}

\end{document}
